% Kunci jawaban Bab 24--32. Berkas ini di-\input langsung di dalam document.
% Setiap angka diuji oleh skrip assert di luar repo (lihat laporan).

\section*{Bab 24: Pemrograman Linear dan Integer}
\begin{enumerate}[label=\textbf{24.\arabic*}]
  \item Maksimumkan $30x_{1}+20x_{2}$ dengan $2x_{1}+x_{2}\le10$, $x_{1}+2x_{2}\le8$,
        $x\ge0$. Verteks $(0;0)$, $(5;0)$, $(0;4)$, $(4;2)$ memberi laba
        $0;\,150;\,80;\,160$. Optimum bersambung $(4;2)$ berlaba $160$; enumerasi bulat
        juga memberi $160$, jadi di sini pembulatan ke bawah kebetulan tepat dan
        celahnya $0$.

  \item Untuk $x,y\in P$ dan $\lambda\in[0,1]$:
        $A(\lambda x+(1-\lambda)y)=\lambda Ax+(1-\lambda)Ay\le\lambda b+(1-\lambda)b=b$,
        memakai $A$ linear dan $\lambda\ge0$. Pada himpunan bulat argumen itu runtuh:
        gabungan konveks dua titik bulat tidak bulat, misalnya $0$ dan $1$ layak
        sementara titik tengahnya bukan anggota $\mathbb{Z}^{n}$.

  \item Laba $x_{1}+2x_{2}$ di verteks: $0$; $4$; $6$; $6$. Optimum $(0;3)$ dan
        $(3;1{,}5)$ sama-sama $6$, dan seluruh ruas $x_{1}+2x_{2}=6$ optimal sebab
        sasarannya konstan di sana; titik tengah $(1{,}5;\,2{,}25)$ bukan verteks.
        Dualnya $y=(0;1)$: harga jam $1$, kayu lapis $0$.

  \item Tabel awal basis $s_{1},s_{2},s_{3}=(4;3;3)$. Pivot pertama: $x_{1}$ masuk
        (biaya tereduksi $3$), nisbah $4$ dan $3$, jadi $s_{2}$ keluar; nilai $9$.
        Pivot kedua: $x_{2}$ masuk, $s_{1}$ keluar (nisbah $1$), optimum $(3;1)$
        dengan $z=11$. Harga dual dari tabel akhir: $(2;1;0)$, dan
        $4(2)+3(1)+3(0)=11$.

  \item Dualnya meminimumkan $10y_{1}+8y_{2}$ dengan $2y_{1}+y_{2}\ge30$,
        $y_{1}+2y_{2}\ge20$, $y\ge0$; penyelesaian $y=(40/3;\,10/3)$ bernilai
        $160$, sama dengan primal. Artinya: satu kg tepung tambahan bernilai
        $13{,}333$ dan satu jam oven tambahan $3{,}333$.

  \item Salah. Pembulatan ke bawah tidak selalu layak begitu ada kendala
        kesamaan atau bertanda $\ge$. Contoh: maksimumkan $x_{1}+x_{2}$ dengan
        $2x_{1}+2x_{2}=5$; jawaban bersambung $(2{,}5;\,0)$ dibulatkan menjadi
        $(2;0)$, yang memberi $4\ne5$ dan karena itu tidak layak.

  \item Benar untuk arah majunya, dan hanya pada pasangan optimal. Kelengkapan
        komplementer $y_{i}(b_{i}-\mathbf{a}_{i}^{\top}\mathbf{x})=0$ memberi:
        $y_{i}>0$ memaksa kendala itu aktif. Sebaliknya tidak berlaku; pada titik
        merosot ada kendala aktif dengan $y_{i}=0$, jadi harga nol tidak berarti
        kendala itu longgar.

  \item Kerapatan $2{,}0$; $1{,}75$; $1{,}667$; $1{,}5$. Barang pertama penuh
        ($5$ kg, $10$), lalu $3/4$ barang kedua, total berat $8$: batas LP
        $15{,}25$. Enumerasi memberi $15$ (barang $1$ dan $3$). Celah $0{,}25$,
        yaitu $1{,}639$ persen dari batas.

  \item 
Kapasitas $7$: barang pertama penuh, lalu $1/2$ barang kedua, batas LP
        $13{,}5$; jawaban bulat terbaik $13$ (barang $1$ dan $4$, berat $7$). Celah
        $0{,}5$, yaitu $3{,}704$ persen. Celah berubah karena kapasitas $7$ membuang
        kombinasi $(1,3)$ yang beratnya $5+3=8$ dan tadinya muat, sehingga jawaban
        bulat terbaik pindah ke barang $1$ dan $4$; pecahan optimum dan bulat
        optimum bergerak tidak sebanding.

  \item $x_{v}=1/2$ untuk semua $v$ layak ($x_{u}+x_{v}=1$) dan memberi $n/2$;
        setiap penutup harus memuat $n-1$ simpul, sebab meninggalkan dua simpul
        berarti rusuk di antaranya tak tertutup. Nisbah LP terhadap bulat:
        $0{,}6667$; $0{,}625$; $0{,}5556$ untuk $n=4,5,10$. Nisbahnya menuju
        $1/2$, jadi batas relaksasi bisa meleset faktor dua dan pembulatan bukan
        cara umum.
\end{enumerate}

\section*{Bab 25: Jaringan Saraf dan Turunan Balik}
\begin{enumerate}[label=\textbf{25.\arabic*}]
  \item Batasnya $\lfloor1/\gamma^{2}\rfloor$: untuk $\gamma=0{,}2$ dapat $25$,
        untuk $\gamma=0{,}05$ dapat $400$. Batas itu hanya menghitung kesalahan,
        dan tiap kesalahan memakan margin, karena itu $n$ tidak muncul. Untuk
        $\gamma$ kecil batasnya meledak seperti $1/\gamma^{2}$ sehingga tak berguna.

  \item Tulis $\mathbf{W}=\mathbf{W}_{2}\mathbf{W}_{1}$ dan
        $\mathbf{b}=\mathbf{W}_{2}\mathbf{b}_{1}+\mathbf{b}_{2}$; maka keluarannya
        $\mathbf{W}\mathbf{x}+\mathbf{b}$, yaitu satu ambang linear. Menambah lapisan
        identitas berikutnya hanya mengalikan matriks, jadi kapasitasnya tidak
        bertambah sedikit pun.

  \item Dengan $\mathbf{x}=(0{,}5;-1;\,0{,}25;2)$ dan $y=0$: $\mathbf{z}_{1}=(2{,}9125;\,
        0{,}375)$, $a=0{,}8961350845$, $\bar{z}_{2}=0{,}8961350845$, sehingga
        $\bar{\mathbf{z}}_{1}=(0{,}7169081;\,-0{,}5376810)$ dan
        \[
          \nabla_{\mathbf{W}_{1}}L=
          \begin{pmatrix}
            0{,}3584540 & -0{,}7169081 & 0{,}1792270 & 1{,}4338161\\
            -0{,}2688405 & 0{,}5376810 & -0{,}1344203 & -1{,}0753621
          \end{pmatrix},
        \]
        yaitu baris kedua sama dengan $-0{,}75$ kali baris pertama. Butir ini
        ambigu: kedua gerbang terbuka, jadi tidak ada baris bergradien nol.

  \item Belum cukup bukti. Selisih hingga hanya memeriksa tiga belas skalar di satu
        titik, dan pada $h=10^{-6}$ galat $10^{-10}$ memang yang diharapkan; bug
        bersama pada kedua hitungan tidak akan terlihat. Perlu: beberapa titik,
        kedua keadaan gerbang, kasus $y=0$ dan $y=1$, gradien terhadap masukan, dan
        pembanding turunan otomatis.

  \item Salah. Teorema itu tentang kepadatan fungsi kontinu pada domain
        kompak: ada parameter yang menghampiri, bukan parameter itu terjangkau
        optimasi. Dua hal yang tidak dinyatakan: banyaknya unit yang diperlukan, dan
        galat pada data uji (generalisasi), bukan hanya pada fungsi sasaran.

  \item Pada $t=1$, $\hat m=g$ dan $\hat v=g^{2}$, sehingga langkahnya
        $\eta g/(|g|+\epsilon)=-0{,}01$ dan $b_{2}$ menjadi $0{,}06$. Dengan gradien
        $100$ kali lebih besar langkahnya tetap $-0{,}01$, sebab Adam hanya
        memakai nisbah terhadap skala gradien; $b_{2}$ tetap $0{,}06$.

  \item 
Keluaran $w_{2}\max(\mathbf{W}_{1}\mathbf{x}+\mathbf{b}_{1},0)$ adalah
        jumlah dua fungsi linear yang masing-masing aktif di satu separuh bidang,
        jadi linear sepotong-sepotong; dua garis $z_{1}[i]=0$ membagi bidang menjadi
        paling banyak empat sel. Untuk $k$ unit dalam susunan umum, $k$ garis
        membagi bidang menjadi $1+k+k(k-1)/2$ daerah, yaitu $2,4,7$ untuk $k=1,2,3$
        (bukan $8$: garis ketiga hanya menambah tiga daerah, bukan empat).

  \item $\nabla_{\mathbf{x}}L=(-0{,}0868370;\,0{,}1328096;\,-0{,}0970532;\,
        -0{,}1430257)$, maka $\boldsymbol{\delta}=0{,}05\,\mathrm{sign}(\nabla)
        =(-0{,}05;0{,}05;-0{,}05;-0{,}05)$. Kerugian di $\mathbf{x}+\boldsymbol{\delta}$
        ialah $0{,}3186800$, sedangkan ramalan linear
        $0{,}2949123+0{,}0229863=0{,}3178985$; selisih $7{,}81\times10^{-4}$ dari
        kelengkungan.

  \item 
Ramalan: blok lab kedua memakai ReLU, dan ReLU hanya jenuh di sisi negatif.
        Tabel saturasi ($8$ unit, $\mathbf{W}_{1}$ He acak, $\mathbf{b}_{2}=0$):
        \begin{center}
        \begin{tabular}{lrrrr}
        & rentang $Z_{1}$ (awal) & unit mati & $L$ akhir & akurasi \\ \hline
        $\mathbf{b}_{1}=0$   & $[-2{,}43;\,2{,}40]$   & $0/8$ & $0{,}0163$ & $1{,}0000$ \\
        $\mathbf{b}_{1}=+10$ & $[7{,}57;\,12{,}41]$   & $0/8$ & $0{,}3225$ & $0{,}8375$ \\
        $\mathbf{b}_{1}=-10$ & $[-12{,}43;\,-7{,}60]$ & $8/8$ & $0{,}6909$ & $0{,}5333$
        \end{tabular}
        \end{center}
        Jadi $\mathbf{b}_{1}=+10$ memang tidak mematikan unit tersembunyi
        ($Z_{1}$ jauh di atas nol diteruskan apa adanya), tetapi tidak gratis:
        $H\cdot\mathbf{w}_{2}$ menjadi sangat besar, sigmoid keluaran menempel di
        $0$ atau $1$, dan $L$ mulai dari $4{,}04$; sesudah $3000$ langkah ia
        berhenti di $L=0{,}3225$ dengan akurasi $0{,}8375$, lebih buruk daripada
        $\mathbf{b}_{1}=0$. Untuk $\mathbf{b}_{1}=-10$ seluruh unit mati sejak awal
        ($Z_{1}$ di bawah nol, turunan tepat $0$, bukan sekadar kecil),
        $\partial L/\partial b_{1}=\mathbf{0}$ tepat, dan latihan berhenti di
        $L=0{,}6909$ -- entropi label -- dengan akurasi kelas mayoritas $0{,}5333$.
  \end{enumerate}

\section*{Bab 26: Generalisasi, Regularisasi, dan Pelatihan}
\begin{enumerate}[label=\textbf{26.\arabic*}]
  \item Belum cukup bukti, dan yang keliru ialah menyamakan galat latih dengan
        galat generalisasi. Galat latih nol hanya berarti data latih terhafal. Angka
        tambahan yang diperlukan: galat validasi (atau uji) pada data yang tidak
        dipakai melatih, beserta lebar selangnya.

  \item Dengan $\hat f=\hat f_{D}$ dan $\bar f=\mathbb{E}_{D}\hat f$:
        $\mathbb{E}(y-\hat f)^{2}=\sigma^{2}+(\bar f-f)^{2}+\operatorname{Var}(\hat f)$.
        Suku silang lenyap karena $\mathbb{E}[\varepsilon]=0$ dan $\varepsilon$
        bebas dari $D$, sehingga $\mathbb{E}[\varepsilon(\hat f-f)]=0$; kebebasan
        itu dipakai tepat di langkah ini, yaitu pada $\varepsilon$, bukan pada
        $\hat f$. Suku silang kedua lenyap dengan alasan berbeda:
        $\mathbb{E}[\hat f-\bar f]=0$ menurut definisi $\bar f$, jadi
        $2(f-\bar f)\,\mathbb{E}[\hat f-\bar f]=0$ tanpa memerlukan kebebasan apa pun.

  \item Dari $\frac{d}{dw}\bigl[\frac12\sum(y_{i}-wx_{i})^{2}\bigr]=0$ dapat
        $w(\lambda)=\sum x_{i}y_{i}/(\sum x_{i}^{2}+\lambda)=64/(30+\lambda)$.
        Turunannya $-64/(30+\lambda)^{2}<0$ tetapi tak pernah nol, jadi $w$ mengecil
        tanpa menyentuh nol. Pada L1, subgradien di $w=0$ memuat nol bila
        $\lambda\ge|\sum x_{i}y_{i}|=64$; jadi $w=0$ tepat pada $\lambda=64$.

  \item $\operatorname{MSE}(c)=(c-1)^{2}\mu^{2}+c^{2}\sigma^{2}/n$: untuk
        $c=0{,}95$ dapat $0{,}01+0{,}0361=0{,}0461$; untuk $c=1{,}05$ dapat
        $0{,}01+0{,}0441=0{,}0541$; untuk $c=1$ dapat $0{,}04$. Simpangan sama besar
        tidak sama rugi karena suku variansi tumbuh seperti $c^{2}$, sehingga
        melebihi jatah lebih mahal daripada menyusut.

  \item Perbaikan: angka validasi silang itu taksiran galat prosedur pemilihan,
        bukan galat model terpilih di lapangan, dan ia bias optimistis karena
        $\lambda$ dipilih pada lipatan yang sama. Perlu dua hal: validasi silang
        bersarang (atau kumpulan uji yang tak pernah dilihat saat memilih), dan
        pengulangan dengan beberapa partisi beserta selangnya.

  \item Untuk $a<0$: dengan $u<0$ turunannya $u-a-t$, nol di $u=a+t$, yang sah
        hanya bila $a<-t$. Untuk $u=0$ subgradiennya $[-a-t,\,-a+t]$, memuat nol
        bila $|a|\le t$, jadi $u=0$ optimal. Simetris karena sasaran
        $\frac12(u-a)^{2}+t|u|$ tidak berubah bila $u$ dan $a$ dua-duanya diganti
        tanda.

  \item Waktu paruhnya $\ln2/(\eta\lambda)=0{,}693147/0{,}002=346{,}57$ langkah;
        seperempat nilai awal berarti dua paruh, $693{,}15$ langkah, dan
        $(1-\eta\lambda)^{693}=0{,}24973$. Simulasi numpy dengan gradien nol
        memberi angka yang sama, sebab $w_{k}=w_{0}(1-\eta\lambda)^{k}$.

  \item Jurangnya $0{,}2406-0{,}1500=0{,}0906$; pada $\lambda=0$ jurangnya
        $0{,}3262-0=0{,}3262$, pada $\lambda=10$ hanya $0{,}3395-0{,}3333=0{,}0062$.
        Jurang kecil tidak cukup: pada $\lambda=10$ jurangnya terkecil justru karena
        keduanya sama-sama buruk, dekat galat kelas mayoritas $0{,}3419$.

  \item Dengan $m$ hipotesis, $P(\exists i:|\hat R_{i}-R_{i}|>\epsilon)\le
        2m e^{-2n\epsilon^{2}}$, jadi $\epsilon=\sqrt{(\ln m+\ln(2/\delta))/(2n)}$.
        Untuk $m=2^{10}$, $n=200$, $\delta=0{,}01$: $\epsilon=0{,}1749$. Separuhnya
        menuntut $n$ empat kali, yaitu $800$.
\end{enumerate}

\section*{Bab 27: Evaluasi Model yang Jujur}
\begin{enumerate}[label=\textbf{27.\arabic*}]
  \item Belum cukup bukti. Selisih $0{,}4$ poin pada satu kumpulan uji belum tentu
        di luar derau; yang keliru ialah memperlakukan dua proporsi berpasangan
        sebagai bebas. Dua angka tambahan: ukuran kumpulan uji, dan cacah
        ketidaksepakatan berpasangan (untuk uji McNemar) atau lebar selang.

  \item Untuk label keras $s,y\in\{0,1\}$ berlaku $(s-y)^{2}=1$ tepat pada
        kesalahan dan $0$ pada kebenaran, jadi rata-ratanya
        $(\mathrm{FP}+\mathrm{FN})/n$. Untuk skor berupa peluang, $(s-y)^{2}$
        mengambil nilai antara, sehingga rata-ratanya bukan lagi cacah kesalahan
        melainkan penalti kalibrasi.

  \item $\mathbb{E}[(s-y)^{2}]=s^{2}-2s\pi+\pi$ dengan $\pi=\mathbb{E}[y]$.
        Turunannya $2s-2\pi$, nol di $s=\pi$; nilai minimumnya
        $\pi^{2}-2\pi^{2}+\pi=\pi(1-\pi)$, dan turunan keduanya $2>0$ memastikan
        itu minimum.

  \item Presisi $180/225=0{,}8$; recall $180/200=0{,}9$; $F_{1}=0{,}847059$;
        akurasi $1935/2000=0{,}9675$; spesifisitas $1755/1800=0{,}975$; akurasi
        seimbang $(0{,}9+0{,}975)/2=0{,}9375$. \texttt{sklearn.metrics} memberi
        angka yang sama sampai enam desimal.

  \item Wilson $95/100$: $[0{,}8882;\,0{,}9785]$, lebar $0{,}0902$; untuk
        $950/1000$: $[0{,}9347;\,0{,}9619]$, lebar $0{,}0272$, yaitu $3{,}318$
        kali lebih sempit meski $n$ sepuluh kali lebih besar. Lebar bukan ukuran
        mutu model: ia bergantung pada $n$ dan pada data, dan kedua percobaan pun
        tak berpasangan.

  \item 
Dari $25$ pasangan, $19$ benar: positif $0{,}95$ mengalahkan lima negatif,
        $0{,}88$ juga lima (ia di atas $0{,}75$), $0{,}70$ empat, $0{,}60$ tiga,
        $0{,}40$ dua. Jadi $\mathrm{AUC}=19/25=0{,}76$; \texttt{roc\_auc\_score}
        memberi $0{,}76$.

  \item Presisi $=0{,}8\pi/(0{,}8\pi+0{,}01(1-\pi))$: pada $\pi=0{,}002$ dapat
        $0{,}138169$, pada $\pi=0{,}2$ dapat $0{,}952381$. Syarat presisi
        $\ge0{,}5$ memberi $0{,}4\pi\ge0{,}005-0{,}005\pi$, yaitu
        $\pi\ge0{,}005/0{,}405=0{,}012346$.

  \item Bagi menurut host, bukan rekaman: $60/20/20$ menjadi $180$ host latih,
        $60$ validasi, $60$ uji, masing-masing sekitar $20$ rekaman. Prapemrosesan
        dihitung di dalam lipatan latih saja. Pemilihan $40$ setelan lewat validasi
        silang bersarang di bagian latih. Uji disentuh satu kali, di akhir.

  \item Salah. AUC mengukur peringkat, bukan kegunaan pada ambang tertentu. Dengan
        sensitivitas $0{,}8$ dan spesifisitas $0{,}99$, pada kelaziman $0{,}002$
        presisinya hanya $0{,}138169$: dari sepuluh alarm, $8{,}62$ di antaranya
        salah, jadi sekitar sembilan. Ambang
        yang dipilih tetap harus diperiksa presisi dan recall-nya.
\end{enumerate}

\section*{Bab 28: Metode Formal dengan Lean 4}
\begin{enumerate}[label=\textbf{28.\arabic*}]
  \item Salah. Berkas yang diterima berarti setiap teorema punya bukti yang
        diperiksa, tetapi \texttt{sorry} lolos dengan peringatan, dan
        \texttt{axiom} maupun \texttt{native\_decide} memperluas dasar percaya.
        Perintah pemeriksa: \texttt{\#print axioms} pada setiap teorema yang
        diklaim.

  \item 
Lemma bantu: untuk setiap $k$ berlaku $k(2k)+2(2k+1)=(k+1)\cdot2(k+1)$.
        Tujuan itu memuat $k\cdot k$, jadi bukan tujuan linear: \texttt{omega} tidak
        menutupnya sendirian dan perlu penulisan ulang lebih dulu. Lalu
        \texttt{induction n with}: basis \texttt{rfl}; langkah \texttt{succ k ih}
        memakai \texttt{rw [oddSum, Nat.mul\_add, ih]}, yang menyisakan tepat lemma
        bantu itu, lalu \texttt{exact bantu k}. Lemma bantu sendiri ditutup
        \texttt{simp only [Nat.mul\_add, Nat.add\_mul, Nat.mul\_one, Nat.one\_mul]}
        diikuti \texttt{omega}: \texttt{omega} tidak bisa menutup tujuan
        $k\cdot(2k)+2(2k+1)=(k+1)\cdot2(k+1)$ sendirian, karena tujuan itu memuat
        $k\cdot k$ dan bukan tujuan linear (dicoba: \texttt{omega could not prove
        the goal}).

  \item Informal: untuk $n=0$ kedua sisi $0$; bila $\sum_{i<n}2i=n(n-1)$, maka
        menambah $2n$ memberi $n(n-1)+2n=(n+1)n$. Versi Lean mendefinisikan
        \texttt{evenSum : Nat -> Nat} rekursif (\texttt{evenSum 0 = 0},
        \texttt{evenSum (n+1) = evenSum n + 2*n}) dan menyatakan bentuk tanpa
        pengurangan terpotong, $E(n)+n=n\cdot n$. Langkahnya: \texttt{induction},
        basis \texttt{rfl}, lalu \texttt{rw [evenSum]}. Perhatikan
        \texttt{rw [evenSum, ih]} langsung gagal (\texttt{did not find instance of
        the pattern evenSum k + k}), sebab tujuan sesudahnya berbunyi
        $E(k)+2k+(k+1)=(k+1)(k+1)$. Selipkan dahulu penataan ulang
        \texttt{have h : evenSum k + 2 * k + (k + 1) = (evenSum k + k) + (2 * k + 1) := by omega},
        lanjutkan \texttt{rw [h, ih]}, dan tutup dengan
        \texttt{simp only [Nat.mul\_add, Nat.add\_mul, Nat.mul\_one, Nat.one\_mul]}
        lalu \texttt{omega}.

  \item Jumlah $37$ bilangan ganjil pertama $=37^{2}=1369$; jumlah $250$ pertama
        $=250^{2}=62500$. Skrip: \texttt{sum(2*i+1 for i in range(n))} memberi
        kedua angka itu, dan keduanya kuadrat sempurna.

  \item        Pemeriksaan langsung untuk semua $n\le1000$ menuntut $1001$ kasus, yaitu
        $n=0$ sampai $n=1000$; batas $n\ge1$ akan memberi $1000$, tetapi soal tidak
        menyebutnya dan basis \texttt{zero} di atas memang memeriksa $n=0$. Bukti
        induksi membangkitkan dua tujuan saja, \texttt{zero} dan \texttt{succ}, lalu
        menyimpulkan seluruh $n$; jadi $1001$ lawan $2$.

  \item \texttt{Nat.add} didefinisikan dengan rekursi pada argumen kedua, jadi
        $n+0$ menyusut langsung ke $n$ dan \texttt{rfl} menutupnya. Pada $0+n$
        rekursinya tersangkut pada $n$ yang belum diketahui, sehingga tidak ada
        yang bisa disusutkan dan diperlukan \texttt{induction n}.

  \item 
Lemah: untuk $x$ tetap, $\exists\delta,\ |\delta|\le1\wedge
        f(x+\delta)=f(x)$; saksi $\delta=0$ membuat kedua sisi sama, \texttt{rfl}
        menutup. Kuat: $\forall x\,\forall\delta,\ |\delta|\le1\to
        f(x+\delta)=f(x)$, dan itu salah untuk $f$ identitas dengan $\delta=1$.

  \item 
\texttt{decide}: tanpa aksioma, murni perhitungan inti. \texttt{omega}:
        \texttt{propext} dan \texttt{Quot.sound}, karena ia membuktikan tujuan
        dengan kontradiksi di dalam bilangan bulat; \texttt{Classical.choice} tidak
        muncul, dan bab ini mencetak daftar yang sama. \texttt{sorry}:
        \texttt{sorryAx}, aksioma yang ditambahkan khusus untuk menutup tujuan yang
        dilewati.

  \item Kesalahan pertama: \texttt{delta} sudah bebas di konteks lalu diikat lagi
        oleh \texttt{exists}, sehingga yang menentukan nilainya pembaca pernyataan,
        bukan penyerang; ketahanan tak tersirat. Kedua: tidak ada kuantor atas $x$,
        jadi $x$ pun bebas. Versi benar:
        $\forall x,\ \exists\delta,\ \lVert\delta\rVert\le1\wedge f(x+\delta)=f(x)$.

  \item \texttt{intro h} memberi $h:P\wedge\neg P$. Komponennya \texttt{h.1} bertipe
        $P$ dan \texttt{h.2} bertipe $\neg P$; lalu \texttt{exact h.2 h.1} menutup
        tujuan, sebab $\neg P$ adalah $P\to$ salah.
\end{enumerate}

\section*{Bab 29: Adversarial Machine Learning: Serangan dan Pertahanan}
\begin{enumerate}[label=\textbf{29.\arabic*}]
  \item $\mathbf{s}_{A}=1$, $\mathbf{s}_{B}=3$, jadi marginnya sudah $-2$ (contoh
        condong ke $B$); $\mathbf{d}=(1;-2;1)$ dengan $\lVert\mathbf{d}\rVert_{1}=4$
        dan $\lVert\mathbf{d}\rVert_{2}=\sqrt6\approx2{,}449$. Maka
        $\varepsilon^{*}_{\infty}=2/4=0{,}5$ dan
        $\varepsilon^{*}_{2}=2/2{,}449=0{,}8165$; rasionya $\sqrt{8/3}=1{,}633$.

  \item Turunan sasaran terhadap $\boldsymbol{\delta}$ ialah $\mathbf{w}$, jadi
        maksimumnya $\varepsilon\lVert\mathbf{w}\rVert_{2}$, dicapai tunggal di
        $\boldsymbol{\delta}=\varepsilon\mathbf{w}/\lVert\mathbf{w}\rVert_{2}$
        (ketaksamaan Cauchy--Schwarz, dengan kesamaan bila sejajar). Pada
        $\ell_\infty$ maksimumnya $\varepsilon\lVert\mathbf{w}\rVert_{1}$ di
        $\varepsilon\,\mathrm{sign}(\mathbf{w})$, sebuah verteks; bentuknya beda
        karena norma dualnya beda.

  \item Galat $\mathbf{d}^{\top}\boldsymbol{\delta}$ linear pada himpunan cembung
        kompak, jadi maksimumnya di titik ekstrem. Titik ekstrem bola
        $\ell_\infty$ berbentuk $\pm\varepsilon$ di tiap koordinat, dan FGSM memilih
        tepat titik itu; karena itu lintasan PGD, yang selalu berada di dalam bola,
        tak dapat melampauinya.

  \item Salah. Pada bab ini gangguan acak bersign di $\varepsilon=0{,}15$ hanya
        menurunkan akurasi ke $95{,}48$ persen, sedangkan PGD 40 langkah
        menjatuhkannya ke $10{,}72$ persen. Ketahanan terhadap derau acak tidak
        mengatakan apa pun tentang serangan yang memilih arah.

  \item 
Keberhasilan tak bersasaran (PGD 40 langkah, akurasi bersih sebagai acuan)
        untuk $A$: $14{,}02$, $53{,}16$, $89{,}06$ persen pada $\varepsilon=
        0{,}05;0{,}10;0{,}15$. Untuk $D$: $5{,}38$, $15{,}28$, $35{,}06$ persen.
        Selisihnya justru melebar pada rentang itu ($8{,}64$, $37{,}88$, $54{,}00$
        poin persen), sebab $A$ sudah mendekati batas $100$ persen sedangkan $D$
        masih jauh; pengecilan baru terjadi di luar rentang, pada
        $\varepsilon\ge0{,}20$, ketika kedua model sama-sama kehabisan margin.

  \item Cukup satu koordinat: proyeksi ke selang adalah pemotongan, dan
        $[x_{j}-\varepsilon,x_{j}+\varepsilon]\cap[0,1]$ sama dengan
        $[\max(0,x_{j}-\varepsilon),\ \min(1,x_{j}+\varepsilon)]$. Memotong ke kotak
        $\varepsilon$ lalu ke $[0,1]$, atau sebaliknya, menghasilkan selang yang
        sama sebab irisan selang bersifat komutatif; hasil kali Cartesiusnya memberi
        titik yang sama.

  \item $D$ kehilangan $97{,}99-96{,}48=1{,}51$ poin akurasi bersih untuk
        memperoleh $62{,}65-10{,}72=51{,}93$ poin ketahanan, jadi harganya
        $1{,}51/51{,}93=0{,}0291$ poin bersih per poin ketahanan. Rasio itu tidak
        lengkap: ia bergantung pada satu $\varepsilon$, pada satu penyerang yang
        bukan batas atas, dan mengabaikan akurasi bersih model pembanding.

  \item Protokol: serang dengan PGD lebih banyak langkah dan beberapa mulai acak,
        dengan EOT atas keacakan pertahanan, lalu ulangi dengan serangan tanpa
        gradien (tiruan/transfer, penaksiran beda hingga). Laporkan kurva akurasi
        terhadap $\varepsilon$, bukan satu titik. Gejala termudah: keberhasilan
        melonjak begitu gradien diganti BPDA atau EOT, sementara serangan naif
        melaporkan nol.

  \item Kotak putih $89{,}06$ persen, transfer $1-0{,}2379/0{,}9849=75{,}85$
        persen, selisih $13{,}21$ poin. Itu bukan bukti $B$ lebih aman: serangan
        yang dibuat khusus untuk $B$ akan lebih kuat, dan di bab ini selisih kotak
        putih terhadap transfer memang menyusut kalau penyerangnya disesuaikan.
        Butir ini ambigu: keluaran lab tidak memuat baris PGD untuk model $B$.

  \item Tak bersasaran cukup mengubah label; bersasaran harus mencapai kelas
        tertentu, jadi himpunan sasaran yang harus dipenuhi jauh lebih kecil. Pada
        $\varepsilon=0{,}15$ sasaran tercapai $33{,}00$ persen, sedangkan tak
        bersasaran $89{,}06$ persen; dengan $\varepsilon=0{,}30$ sasaran tercapai
        $99{,}50$ persen. Jadi kesulitannya soal anggaran, bukan jenis serangan.
\end{enumerate}

\section*{Bab 30: Batas Serangan: Gradien Dikaburkan dan Evasi NIDS}
\begin{enumerate}[label=\textbf{30.\arabic*}]
  \item Tiga pemeriksaan: (a) ulangi dengan serangan adaptif (BPDA, EOT), (b)
        laporkan kurva akurasi terhadap $\varepsilon$, bukan satu titik, (c) periksa
        mutu gradien (kosinus terhadap beda hingga atau terhadap arah bobot). Bila
        yang terjadi penyamaran gradien: keberhasilan hanya runtuh untuk serangan
        berbasis gradien, sementara pencarian acak dan transfer tetap berhasil, dan
        kurvanya datar.

  \item Untuk $f(x)=w\cdot x+b$ gradiennya konstan, Hessiannya nol, jadi konstanta
        Lipschitz $L=0$ dan batas $L\delta$ memberi galat nol: BPDA tepat. Pada dua
        lapis ReLU, $f$ linear sepotong-sepotong sehingga gradiennya melompat di
        patahan dan hipotesis gradien Lipschitz gagal; batas itu lalu tak berarti.

  \item 
Dengan program: label berbalik sampai $q=0{,}6$ (skor $-1{,}2$) dan gagal
        pada $q\ge0{,}70$, ketika semua koordinat membulat ke nol sehingga skornya
        $0$, yaitu seri, bukan pembalikan. Jadi
        $q^{*}=2\max_{j}|x_{j}-\varepsilon|=2(\varepsilon+|\min_{j}x_{j}|)=0{,}70$.
        Batas aman analitis: galat pembulatan paling besar $4q$, jadi serangan hanya
        dijamin berhasil bila $q<(8\varepsilon-m)/4=0{,}345$; pada $q=0{,}345$ label
        masih berbalik, sehingga arah ketaksamaannya penting.

  \item Peluang menolak $=P(\text{lebih dari }10\text{ dari }20\text{ hilang})
        =0{,}411901$, jadi yang dijawab $588$ dari $1000$ kueri: $588/20=29$
        langkah PGD dengan $m=20$. Dengan $m=1$ peluang menolak $0{,}5$, jadi $500$
        langkah. Yang dibayar: gradien makin berderau, arahnya makin jauh dari benar.

  \item Yang dapat diramalkan tanpa menjalankan: akurasi bersih hampir sama,
        keberhasilan naif tetap mendekati nol, dan keberhasilan BPDA mendekati garis
        dasar untuk $q=0$ dan $q=0{,}05$ tetapi turun untuk $q$ besar. Kosinusnya
        tidak lagi mendekati satu karena gerbang ReLU menambah patahan di atas
        pembulatan, sehingga aproksimasi BPDA hanya menghampiri arah. Angka ketiga
        nilai $q$ harus diambil dari eksekusi, bukan dikarang.

  \item Bendera biner dan nomor protokol dianggap tak dapat diubah penyerang:
        $7+3=10$ fitur terkunci, sisa $40$ fitur yang benar-benar bisa digeser (dari
        $50$). Gambar $32\times32$ dengan tiga kanal punya $3072$ nilai, jadi
        $3072/40=76{,}8$ kali lebih banyak.

  \item Pelatihan ulang harian mengubah parameter yang dipakai serangan bergradien,
        jadi arah lama cepat basi; serangan buta tidak memakai parameter itu, jadi
        ia tetap bekerja. Menurut lemma EOT yang harus diperiksa adalah harapan atas
        transformasi, dan serangan adaptif harus didefinisikan terhadap pertahanan
        itu sendiri, termasuk jadwal pelatihan ulangnya.

  \item Metrik: keberhasilan serangan adaptif terhadap $\varepsilon$, dengan
        sekurangnya $5$ seed dan selang. Bentuk kurva yang menerima (a): kurva
        tergeser ke kanan atau ke bawah secara sistematis pada semua $\varepsilon$,
        bukan hanya di satu titik. Bila setelah BPDA dan EOT keberhasilannya kembali
        ke garis dasar, yang rusak hanya pengukuran, jadi (b).
\end{enumerate}

\section*{Bab 31: Kriptografi: Bilangan Bulat Modular sampai Ancaman Kuantum}
\begin{enumerate}[label=\textbf{31.\arabic*}]
  \item Salah. $2^{100}\equiv1\pmod{101}$ hanya mengatakan orde $2$ membagi $100$,
        yaitu $101$ lolos uji Fermat basis $2$; bilangan semu (dan bilangan
        Carmichael) juga lolos. Yang dibuktikan kongruensi itu hanyalah $101\nmid2$
        dan $2^{100}\equiv1$, bukan ketiadaan pembagi.

  \item Bila $\gcd(a,n)=\gcd(b,n)=1$, teorema Bezout memberi $ax+ny=1$ dan
        $bz+nw=1$; hasilkali keduanya memberi $(ab)(xz)+n(\cdot)=1$, jadi $ab$ unit
        dan $U(n)$ tertutup. Bila $a$ bukan unit, $d=\gcd(a,n)>1$ dan
        $b=n/d\ne0$ memenuhi $ab=n(a/d)\equiv0$.

  \item $a^{2}-1=(a-1)(a+1)\equiv0\pmod p$; karena $p$ prima, lemma Euklid memberi
        $p\mid a-1$ atau $p\mid a+1$, jadi $a\equiv\pm1$. Untuk modulus $8$ lemma itu
        gugur sebab $8$ bukan prima: akar kuadrat dari satu ialah $1,3,5,7$, karena
        $3^{2}=5^{2}=7^{2}=49\equiv1$.

  \item Kelas $0\le a<pq$ yang berbagi faktor dengan $pq$ ada $q$ kelipatan $p$
        ditambah $p$ kelipatan $q$, dan $0$ terhitung dua kali; jadi
        $\varphi(pq)=pq-p-q+1=(p-1)(q-1)$. Untuk $p=1019$, $q=1031$:
        $\varphi=1018\cdot1030=1\,048\,540$.

  \item Euklid mundur memberi $17^{-1}\equiv2753\pmod{3120}$ (sebab
        $17\cdot2753=46\,801=15\cdot3120+1$), dan
        $65\,537^{-1}\equiv60\,493\pmod{1\,048\,540}$. Pembagiannya $4$ lawan $5$
        langkah, jadi menambah satu digit pada modulus hanya menambah sedikit
        langkah.

  \item Fermat kecil: $a^{100}\equiv1\pmod{101}$ untuk $101\nmid a$. Maka
        $5^{100}\equiv1$, $3^{200}=(3^{100})^{2}\equiv1$, dan
        $2^{1000}=(2^{100})^{10}\equiv1$. Ketiganya diperiksa
        \texttt{pow(a,k,101)} dan cocok.

  \item $A=5^{7}\equiv17$, $B=5^{15}\equiv19$, kunci bersama $19^{7}\equiv15$
        (sama dengan $17^{15}$). Bila penyadap mengganti $A$ dengan $11$, Bob
        menghitung $11^{15}\equiv10$ sedangkan Alice tetap memegang $15$; dua kunci
        berbeda muncul karena pangkat yang dipasangkan bukan lagi pasangan
        Diffie--Hellman yang sah.

  \item Dengan $e=3$ dan $m<n^{1/3}$ berlaku $m^{3}<n$, jadi $c=m^{3}$ tanpa
        pemotongan modulus dan akar pangkat tiga bilangan bulat mengembalikan $m$.
        Untuk $n=1\,050\,589$ batasnya $n^{1/3}\approx101{,}66$, jadi pesan
        terpanjang $m=101$. Bantalan acak membuat $m^{3}>n$ sehingga akar bulat itu
        tak ada.

  \item Dua hal yang harus benar: (i) parameter dan implementasi benar-benar
        memenuhi jaminan keamanan yang diklaim, yaitu reduksi terhadap masalah
        kisi yang masih sulit; (ii) asumsi kekerasan yang mendasarinya masih
        berlaku dan model pembuktiannya mencakup serangan yang benar-benar dipakai.
        Adopsi luas bukan syarat kebenaran klaim: skema yang dipakai di mana pun
        tetap tidak terbukti aman bila asumsinya gugur.
        Kalimat kedua tetap keliru: RSA masih tertanam di sistem lama, dan migrasi
        dilakukan hibrida, jadi butuh bertahun-tahun. (Tanda tangan pascakuantum
        sudah ada: ML-DSA, FIPS 204.)

  \item Sekitar $1{,}1774\sqrt{2^{48}}=1{,}1774\cdot2^{24}$, yaitu $19\,753\,662$
        sampel untuk peluang tabrakan seperdua; untuk $32$ bit
        $1{,}1774\cdot2^{16}=77\,163$. Angka $2^{24}$ dan $2^{16}$ adalah tebakan
        kasar akar kuadrat, bukan nilai tepatnya.
\end{enumerate}

\section*{Bab 32: Privasi Diferensial dan Fondasi Matematis QML}
\begin{enumerate}[label=\textbf{32.\arabic*}]
  \item Salah sebagai kalimat. $\varepsilon=0{,}5$ membatasi nisbah peluang per
        catatan, bukan menghapus kebocoran; ia juga jaminan kasus terburuk. Dua hal
        yang harus benar: mekanismenya memang memenuhi $(\varepsilon,\delta)$ untuk
        relasi ketetanggaan dan sensitivitas yang dipakai, dan seluruh keluaran
        (kueri, sisa anggaran, pascaproses) dihitung dalam anggaran itu.

  \item Kepadatan Laplace berdimensi $d$ adalah hasil kali, jadi rasionya
        $\prod_{j}\exp\bigl((|y_{j}-a'_{j}|-|y_{j}-a_{j}|)/b\bigr)$, dan dengan
        ketaksamaan segitiga tiap faktor paling besar
        $\exp(|a_{j}-a'_{j}|/b)$. Hasilkalinya
        $\exp(\sum_{j}|a_{j}-a'_{j}|/b)\le\exp(\Delta f/b)=e^{\varepsilon}$.
        Argumen satu dimensi diganti hasil kali rasio per koordinat, dan
        $\ell_1$ menggantikan nilai mutlak tunggal.

  \item Ambil $D\sim D'$ yang berbeda pada satu catatan. Catatan itu hanya ada di
        satu dari $D_{1},D_{2}$, jadi hanya satu mekanisme yang masukannya berubah;
        rasionya dibatasi $e^{\varepsilon_{i}}$, sedangkan yang lain berasio $1$.
        Karena itu $\varepsilon=\max_{i}\varepsilon_{i}$ dan
        $\delta=\max_{i}\delta_{i}$, bukan jumlah.

  \item Peluang menjawab ``ya'' ialah $p+(1-p)/2$ bila benar ``ya'' dan $(1-p)/2$
        bila benar ``tidak''; rasionya $(1+p)/(1-p)$. Maka
        $p=(e^{\varepsilon}-1)/(e^{\varepsilon}+1)$, dan untuk $\varepsilon=2$
        didapat $p=0{,}761594$, yaitu $\tanh1$.

  \item Kueri hitung: $b=2$, jadi $P(|Y|>10)=e^{-5}=0{,}006738$. Kueri rata-rata
        dengan $R=1$: $b=10^{-4}/0{,}5=2\times10^{-4}$; galat $10$ orang berarti
        $10^{-3}$ pada skala rata-rata, yang kembali memberi $e^{-5}$, sedangkan
        galat mutlak $10$ pada skala itu memberi $e^{-50000}\approx0$. Butir ini
        ambigu soal satuan galatnya.

  \item $H|0\rangle=(|0\rangle+|1\rangle)/\sqrt2$; $S$ memberi
        $(|0\rangle+i|1\rangle)/\sqrt2$; lalu $H$ memberi
        $\frac{1+i}{2}|0\rangle+\frac{1-i}{2}|1\rangle$. Peluangnya $1/2$ dan $1/2$,
        sama dengan $H|0\rangle$, sebab ukuran basis komputasi hanya melihat modulus
        amplitudo. Yang membedakannya fase relatif, dan gerbangnya $S^{\dagger}$
        (memutar $|1\rangle$ dengan $-i$). $Z$ sesudah $H$ bukan pengganti: $HZH=X$,
        yang menukar $|0\rangle$ dengan $|1\rangle$, bukan memutar fasenya.

  \item $|\Phi^{+}\rangle=(|00\rangle+|11\rangle)/\sqrt2$. Bila qubit kedua diukur
        lebih dulu dan hasilnya $0$, qubit pertama runtuh ke $|0\rangle$, jadi
        peluang hasil $01$ ialah $0$; bila hasilnya $1$, qubit pertama runtuh ke
        $|1\rangle$, jadi hasil $11$ pasti dan peluang $01$ juga $0$. Keduanya
        menunjukkan korelasi sempurna: hasil kedua qubit selalu sama.

  \item $U^{\dagger}U=\operatorname{diag}(1,e^{-i\varphi})
        \operatorname{diag}(1,e^{i\varphi})=I$, jadi uniter untuk setiap
        $\varphi$. Pada basis komputasi $|1\rangle$ hanya mendapat fase, modulus
        amplitudonya tetap, sehingga peluang tak berubah; tetapi fase relatif itu
        mengubah interferensi pada gerbang $H$ berikutnya, sehingga hasilnya bisa
        berbeda.

  \item 
$n=40$: $2^{40}=1{,}0995\times10^{12}$ amplitudo, kali $16$ byte menjadi
        $16$ TiB (bukan GiB: $2^{40}$ byte sudah satu TiB). $n=60$:
        $2^{60}=1{,}1529\times10^{18}$ amplitudo, yaitu $16$ EiB atau
        $1{,}8447\times10^{19}$ byte, jauh melampaui penyimpanan mana pun yang
        tersedia.

  \item Tiga pertanyaan: apakah seluruh amplitudo benar-benar dipakai, apakah
        keluarannya dapat diekstrak dengan jumlah pengukuran yang tidak eksponensial,
        dan apakah ada pembanding klasik pada tugas dan data yang sama. Yang membuat
        saya menerima: ada pemisahan yang terbukti atau penskalaan terukur pada tugas
        itu, jumlah pengukuran polinomial, dan keunggulan atas garis dasar klasik
        yang disetel serius.
\end{enumerate}
